% --- Archon physics formalization source begin ---
% archon:physics
% archon:covers PhyXMiniProblems/problem_phyx_mini_0295.lean
% archon:source-report reports/phyx_mini/problem_phyx_mini_0295.source.json

\chapter{Physics problem phyx\_mini\_0295}
\label{ch:PhyXMiniProblems_problem_phyx_mini_0295}

\paragraph{Problem source.}
\#\# Physical scenario

A nylon guitar string has a linear density of \$7.20\$ g/m and is under a tension of \$150\$ N. The fixed supports are distance \$D = 90.0\$ cm apart. The string is oscillating in the standing wave pattern shown in figure.

\#\# Question

Calculate the frequency of the traveling waves whose superposition gives this standing wave.

\#\# Answer choices

- A: A: 226 Hz

- B: B: 230 Hz

- C: C: 235 Hz

- D: D: 241 Hz

\#\# Figure caption (auxiliary; use the image as primary evidence)

The image is a diagram featuring a sine wave and arrows. Here is a detailed description:

1. **Sine Wave**: 
   - An orange sine wave runs horizontally across the image.
   - The sine wave has a smooth, flowing curvature with alternating solid and dashed segments.

2. **Arrows**:
   - Above the sine wave, there is a horizontal double-headed arrow.
   - The arrow points left and right, indicating a distance or measurement.

3. **Text**:
   - The letter "D" is centered above the sine wave and between the two arrowheads. 
   - This suggests that "D" represents the distance or wavelength between points on the wave.

4. **Dots**:
   - At each end of the sine wave and arrows, there are black dots.

The diagram seems to illustrate a concept related to wave physics, such as wavelength or distance between wave peaks.

\#\# Dataset metadata

Resonance and Harmonics; Physical Model Grounding Reasoning, Multi-Formula Reasoning

\paragraph{Recorded answer/context.}
D: D: 241 Hz

\paragraph{Figure/image path.}
/root/proposal\_for\_physic/science-mango/phyx\_mini\_run/phyx\_data/test\_image/295.png

\paragraph{Formalization target.}
create a compiling Lean file with sorry bodies at `PhyXMiniProblems/problem\_phyx\_mini\_0295.lean`. The Lean declarations must preserve the physical quantities, dimensions or dimensional roles, figure labels, governing-law hypotheses, and final relation expressed by this problem.
Use Mathlib/Physlib names found through LeanExplore where available. If a physics API is missing, introduce faithful local abstractions rather than scalar placeholder aliases.

\begin{theorem}[Physics formalization target]
\label{thm:physics:phyx_mini_0295:target}
\lean{PhyXMiniProblems.ProblemPhyXMini0295.travelingWaveFrequency_matches_recordedAnswer}
\uses{def:physics:phyx-mini-0295:phyxminiproblems-problemphyxmini0295-guitarstringstandingwavesetup, def:physics:phyx-mini-0295:phyxminiproblems-problemphyxmini0295-matchesproblemandprimaryfigure, def:physics:phyx-mini-0295:phyxminiproblems-problemphyxmini0295-haspositivewaveparameters, def:physics:phyx-mini-0295:phyxminiproblems-problemphyxmini0295-satisfiesfixedendstandingwavegeometry, def:physics:phyx-mini-0295:phyxminiproblems-problemphyxmini0295-satisfiestravelingwaverelation, def:physics:phyx-mini-0295:phyxminiproblems-problemphyxmini0295-satisfiestautstringwavespeedlaw, def:physics:phyx-mini-0295:phyxminiproblems-problemphyxmini0295-recordedanswerchoice, def:physics:phyx-mini-0295:phyxminiproblems-problemphyxmini0295-matchesanswerchoice}
The assigned autoformalize agent should translate this physics problem into Lean declarations in the covered file, with theorem and lemma proof bodies written as `by sorry`.
\end{theorem}
\begin{proof}
This is an autoformalization task, not a proof task. Produce statements that can later be proved by the physics prover without weakening the physical model.
\end{proof}

% --- Archon declaration topology begin ---
\section{Declaration topology}
\begin{definition}[Length Quantity]
  \label{def:physics:phyx-mini-0295:phyxminiproblems-problemphyxmini0295-lengthquantity}
  \lean{PhyXMiniProblems.ProblemPhyXMini0295.LengthQuantity}
  A physical length
\end{definition}
\begin{proof}
  This is the declared model definition of Length Quantity.
\end{proof}
\begin{definition}[Frequency Quantity]
  \label{def:physics:phyx-mini-0295:phyxminiproblems-problemphyxmini0295-frequencyquantity}
  \lean{PhyXMiniProblems.ProblemPhyXMini0295.FrequencyQuantity}
  A physical frequency, carrying the inverse-time dimension
\end{definition}
\begin{proof}
  This is the declared model definition of Frequency Quantity.
\end{proof}
\begin{definition}[Linear Mass Density Quantity]
  \label{def:physics:phyx-mini-0295:phyxminiproblems-problemphyxmini0295-linearmassdensityquantity}
  \lean{PhyXMiniProblems.ProblemPhyXMini0295.LinearMassDensityQuantity}
  A physical mass per unit length
\end{definition}
\begin{proof}
  This is the declared model definition of Linear Mass Density Quantity.
\end{proof}
\begin{definition}[Speed Quantity]
  \label{def:physics:phyx-mini-0295:phyxminiproblems-problemphyxmini0295-speedquantity}
  \lean{PhyXMiniProblems.ProblemPhyXMini0295.SpeedQuantity}
  A physical propagation speed
\end{definition}
\begin{proof}
  This is the declared model definition of Speed Quantity.
\end{proof}
\begin{definition}[Force Quantity]
  \label{def:physics:phyx-mini-0295:phyxminiproblems-problemphyxmini0295-forcequantity}
  \lean{PhyXMiniProblems.ProblemPhyXMini0295.ForceQuantity}
  A physical force magnitude, used here for string tension
\end{definition}
\begin{proof}
  This is the declared model definition of Force Quantity.
\end{proof}
\begin{definition}[quantity Readout]
  \label{def:physics:phyx-mini-0295:phyxminiproblems-problemphyxmini0295-quantityreadout}
  \lean{PhyXMiniProblems.ProblemPhyXMini0295.quantityReadout}
  Read a nonnegative, dimension-tagged physical quantity in chosen units
\end{definition}
\begin{proof}
  This is the declared model definition of quantity Readout.
\end{proof}
\begin{definition}[length In Meters]
  \label{def:physics:phyx-mini-0295:phyxminiproblems-problemphyxmini0295-lengthinmeters}
  \lean{PhyXMiniProblems.ProblemPhyXMini0295.lengthInMeters}
  \uses{def:physics:phyx-mini-0295:phyxminiproblems-problemphyxmini0295-lengthquantity, def:physics:phyx-mini-0295:phyxminiproblems-problemphyxmini0295-quantityreadout}
  Metre readout of a physical length
\end{definition}
\begin{proof}
  This is the declared model definition of length In Meters.
\end{proof}
\begin{definition}[frequency In Hertz]
  \label{def:physics:phyx-mini-0295:phyxminiproblems-problemphyxmini0295-frequencyinhertz}
  \lean{PhyXMiniProblems.ProblemPhyXMini0295.frequencyInHertz}
  \uses{def:physics:phyx-mini-0295:phyxminiproblems-problemphyxmini0295-frequencyquantity, def:physics:phyx-mini-0295:phyxminiproblems-problemphyxmini0295-quantityreadout}
  Hertz readout of a physical frequency
\end{definition}
\begin{proof}
  This is the declared model definition of frequency In Hertz.
\end{proof}
\begin{definition}[linear Mass Density In Kilograms Per Meter]
  \label{def:physics:phyx-mini-0295:phyxminiproblems-problemphyxmini0295-linearmassdensityinkilogramspermeter}
  \lean{PhyXMiniProblems.ProblemPhyXMini0295.linearMassDensityInKilogramsPerMeter}
  \uses{def:physics:phyx-mini-0295:phyxminiproblems-problemphyxmini0295-linearmassdensityquantity, def:physics:phyx-mini-0295:phyxminiproblems-problemphyxmini0295-quantityreadout}
  Kilograms-per-metre readout of a physical linear mass density
\end{definition}
\begin{proof}
  This is the declared model definition of linear Mass Density In Kilograms Per Meter.
\end{proof}
\begin{definition}[speed In Meters Per Second]
  \label{def:physics:phyx-mini-0295:phyxminiproblems-problemphyxmini0295-speedinmeterspersecond}
  \lean{PhyXMiniProblems.ProblemPhyXMini0295.speedInMetersPerSecond}
  \uses{def:physics:phyx-mini-0295:phyxminiproblems-problemphyxmini0295-speedquantity, def:physics:phyx-mini-0295:phyxminiproblems-problemphyxmini0295-quantityreadout}
  Metres-per-second readout of a physical speed
\end{definition}
\begin{proof}
  This is the declared model definition of speed In Meters Per Second.
\end{proof}
\begin{definition}[force In Newtons]
  \label{def:physics:phyx-mini-0295:phyxminiproblems-problemphyxmini0295-forceinnewtons}
  \lean{PhyXMiniProblems.ProblemPhyXMini0295.forceInNewtons}
  \uses{def:physics:phyx-mini-0295:phyxminiproblems-problemphyxmini0295-forcequantity, def:physics:phyx-mini-0295:phyxminiproblems-problemphyxmini0295-quantityreadout}
  Newton readout of a force magnitude in coherent SI base units
\end{definition}
\begin{proof}
  This is the declared model definition of force In Newtons.
\end{proof}
\begin{definition}[String Material]
  \label{def:physics:phyx-mini-0295:phyxminiproblems-problemphyxmini0295-stringmaterial}
  \lean{PhyXMiniProblems.ProblemPhyXMini0295.StringMaterial}
  Material classification of the vibrating string
\end{definition}
\begin{proof}
  This is the declared model definition of String Material.
\end{proof}
\begin{definition}[String Endpoint]
  \label{def:physics:phyx-mini-0295:phyxminiproblems-problemphyxmini0295-stringendpoint}
  \lean{PhyXMiniProblems.ProblemPhyXMini0295.StringEndpoint}
  The two supports delimiting the vibrating part of the string
\end{definition}
\begin{proof}
  This is the declared model definition of String Endpoint.
\end{proof}
\begin{definition}[Endpoint Boundary]
  \label{def:physics:phyx-mini-0295:phyxminiproblems-problemphyxmini0295-endpointboundary}
  \lean{PhyXMiniProblems.ProblemPhyXMini0295.EndpointBoundary}
  Mechanical boundary condition at an endpoint of the string
\end{definition}
\begin{proof}
  This is the declared model definition of Endpoint Boundary.
\end{proof}
\begin{definition}[Figure Label]
  \label{def:physics:phyx-mini-0295:phyxminiproblems-problemphyxmini0295-figurelabel}
  \lean{PhyXMiniProblems.ProblemPhyXMini0295.FigureLabel}
  The only text label in the primary figure
\end{definition}
\begin{proof}
  This is the declared model definition of Figure Label.
\end{proof}
\begin{definition}[Figure Span]
  \label{def:physics:phyx-mini-0295:phyxminiproblems-problemphyxmini0295-figurespan}
  \lean{PhyXMiniProblems.ProblemPhyXMini0295.FigureSpan}
  The geometric span to which a figure label can point
\end{definition}
\begin{proof}
  This is the declared model definition of Figure Span.
\end{proof}
\begin{definition}[Guitar String Standing Wave Setup]
  \label{def:physics:phyx-mini-0295:phyxminiproblems-problemphyxmini0295-guitarstringstandingwavesetup}
  \lean{PhyXMiniProblems.ProblemPhyXMini0295.GuitarStringStandingWaveSetup}
  \uses{def:physics:phyx-mini-0295:phyxminiproblems-problemphyxmini0295-lengthquantity, def:physics:phyx-mini-0295:phyxminiproblems-problemphyxmini0295-frequencyquantity, def:physics:phyx-mini-0295:phyxminiproblems-problemphyxmini0295-linearmassdensityquantity, def:physics:phyx-mini-0295:phyxminiproblems-problemphyxmini0295-speedquantity, def:physics:phyx-mini-0295:phyxminiproblems-problemphyxmini0295-forcequantity, def:physics:phyx-mini-0295:phyxminiproblems-problemphyxmini0295-stringmaterial, def:physics:phyx-mini-0295:phyxminiproblems-problemphyxmini0295-stringendpoint, def:physics:phyx-mini-0295:phyxminiproblems-problemphyxmini0295-endpointboundary, def:physics:phyx-mini-0295:phyxminiproblems-problemphyxmini0295-figurelabel, def:physics:phyx-mini-0295:phyxminiproblems-problemphyxmini0295-figurespan}
  Physical quantities and figure-derived roles for the vibrating string. The wavelength, propagation speed, and component traveling-wave frequency are unknown state variables. In particular, this structure assigns none of them the requested numerical answer
\end{definition}
\begin{proof}
  This is the declared model definition of Guitar String Standing Wave Setup.
\end{proof}
\begin{definition}[Matches Problem And Primary Figure]
  \label{def:physics:phyx-mini-0295:phyxminiproblems-problemphyxmini0295-matchesproblemandprimaryfigure}
  \lean{PhyXMiniProblems.ProblemPhyXMini0295.MatchesProblemAndPrimaryFigure}
  \uses{def:physics:phyx-mini-0295:phyxminiproblems-problemphyxmini0295-lengthinmeters, def:physics:phyx-mini-0295:phyxminiproblems-problemphyxmini0295-linearmassdensityinkilogramspermeter, def:physics:phyx-mini-0295:phyxminiproblems-problemphyxmini0295-forceinnewtons, def:physics:phyx-mini-0295:phyxminiproblems-problemphyxmini0295-guitarstringstandingwavesetup}
  The categorical and numerical information supplied by the problem and primary figure. The figure has nodes at its black endpoint dots and three antinodes; the arrow labelled D spans those endpoints. The density conversion is 7.20 g/m = 9/1250 kg/m, and 90.0 cm = 9/10 m. No wavelength, wave speed, or frequency value is included here
\end{definition}
\begin{proof}
  This is the declared model definition of Matches Problem And Primary Figure.
\end{proof}
\begin{definition}[Has Positive Wave Parameters]
  \label{def:physics:phyx-mini-0295:phyxminiproblems-problemphyxmini0295-haspositivewaveparameters}
  \lean{PhyXMiniProblems.ProblemPhyXMini0295.HasPositiveWaveParameters}
  \uses{def:physics:phyx-mini-0295:phyxminiproblems-problemphyxmini0295-lengthinmeters, def:physics:phyx-mini-0295:phyxminiproblems-problemphyxmini0295-frequencyinhertz, def:physics:phyx-mini-0295:phyxminiproblems-problemphyxmini0295-linearmassdensityinkilogramspermeter, def:physics:phyx-mini-0295:phyxminiproblems-problemphyxmini0295-speedinmeterspersecond, def:physics:phyx-mini-0295:phyxminiproblems-problemphyxmini0295-forceinnewtons, def:physics:phyx-mini-0295:phyxminiproblems-problemphyxmini0295-guitarstringstandingwavesetup}
  Positivity assumptions selecting the nondegenerate physical wave state
\end{definition}
\begin{proof}
  This is the declared model definition of Has Positive Wave Parameters.
\end{proof}
\begin{definition}[Satisfies Fixed End Standing Wave Geometry]
  \label{def:physics:phyx-mini-0295:phyxminiproblems-problemphyxmini0295-satisfiesfixedendstandingwavegeometry}
  \lean{PhyXMiniProblems.ProblemPhyXMini0295.SatisfiesFixedEndStandingWaveGeometry}
  \uses{def:physics:phyx-mini-0295:phyxminiproblems-problemphyxmini0295-lengthinmeters, def:physics:phyx-mini-0295:phyxminiproblems-problemphyxmini0295-guitarstringstandingwavesetup}
  The general fixed-end geometry n * lambda = 2 * D, where n is the number of half-wavelength loops. This law does not prescribe the three-loop mode
\end{definition}
\begin{proof}
  This is the declared model definition of Satisfies Fixed End Standing Wave Geometry.
\end{proof}
\begin{definition}[Satisfies Traveling Wave Relation]
  \label{def:physics:phyx-mini-0295:phyxminiproblems-problemphyxmini0295-satisfiestravelingwaverelation}
  \lean{PhyXMiniProblems.ProblemPhyXMini0295.SatisfiesTravelingWaveRelation}
  \uses{def:physics:phyx-mini-0295:phyxminiproblems-problemphyxmini0295-lengthinmeters, def:physics:phyx-mini-0295:phyxminiproblems-problemphyxmini0295-frequencyinhertz, def:physics:phyx-mini-0295:phyxminiproblems-problemphyxmini0295-speedinmeterspersecond, def:physics:phyx-mini-0295:phyxminiproblems-problemphyxmini0295-guitarstringstandingwavesetup}
  The nondispersive traveling-wave relation v = f * lambda
\end{definition}
\begin{proof}
  This is the declared model definition of Satisfies Traveling Wave Relation.
\end{proof}
\begin{definition}[Satisfies Taut String Wave Speed Law]
  \label{def:physics:phyx-mini-0295:phyxminiproblems-problemphyxmini0295-satisfiestautstringwavespeedlaw}
  \lean{PhyXMiniProblems.ProblemPhyXMini0295.SatisfiesTautStringWaveSpeedLaw}
  \uses{def:physics:phyx-mini-0295:phyxminiproblems-problemphyxmini0295-linearmassdensityinkilogramspermeter, def:physics:phyx-mini-0295:phyxminiproblems-problemphyxmini0295-speedinmeterspersecond, def:physics:phyx-mini-0295:phyxminiproblems-problemphyxmini0295-forceinnewtons, def:physics:phyx-mini-0295:phyxminiproblems-problemphyxmini0295-guitarstringstandingwavesetup}
  The transverse-wave law for a uniform taut string, in the homogeneous squared form T = mu * v\textasciicircum{}2 in coherent SI readouts
\end{definition}
\begin{proof}
  This is the declared model definition of Satisfies Taut String Wave Speed Law.
\end{proof}
\begin{lemma}[component Traveling Wave Frequency In Hertz eq exact]
  \label{lem:physics:phyx-mini-0295:phyxminiproblems-problemphyxmini0295-componenttravelingwavefrequencyinhertz-eq-exact}
  \lean{PhyXMiniProblems.ProblemPhyXMini0295.componentTravelingWaveFrequencyInHertz_eq_exact}
  \uses{def:physics:phyx-mini-0295:phyxminiproblems-problemphyxmini0295-frequencyinhertz, def:physics:phyx-mini-0295:phyxminiproblems-problemphyxmini0295-guitarstringstandingwavesetup, def:physics:phyx-mini-0295:phyxminiproblems-problemphyxmini0295-matchesproblemandprimaryfigure, def:physics:phyx-mini-0295:phyxminiproblems-problemphyxmini0295-haspositivewaveparameters, def:physics:phyx-mini-0295:phyxminiproblems-problemphyxmini0295-satisfiesfixedendstandingwavegeometry, def:physics:phyx-mini-0295:phyxminiproblems-problemphyxmini0295-satisfiestravelingwaverelation, def:physics:phyx-mini-0295:phyxminiproblems-problemphyxmini0295-satisfiestautstringwavespeedlaw}
  The three-antinode geometry and taut-string laws determine the unrounded component-wave frequency as 1250 * sqrt 3 / 9 hertz
\end{lemma}
\begin{proof}
  The conclusion follows from the governing relations and figure readouts named in the statement of component Traveling Wave Frequency In Hertz eq exact.
\end{proof}
\begin{definition}[Answer Choice]
  \label{def:physics:phyx-mini-0295:phyxminiproblems-problemphyxmini0295-answerchoice}
  \lean{PhyXMiniProblems.ProblemPhyXMini0295.AnswerChoice}
  Labels of the four displayed answer choices
\end{definition}
\begin{proof}
  This is the declared model definition of Answer Choice.
\end{proof}
\begin{definition}[frequency In Hertz]
  \label{def:physics:phyx-mini-0295:phyxminiproblems-problemphyxmini0295-answerchoice-frequencyinhertz}
  \lean{PhyXMiniProblems.ProblemPhyXMini0295.AnswerChoice.frequencyInHertz}
  \uses{def:physics:phyx-mini-0295:phyxminiproblems-problemphyxmini0295-frequencyinhertz, def:physics:phyx-mini-0295:phyxminiproblems-problemphyxmini0295-answerchoice}
  Frequency in hertz printed beside each answer label
\end{definition}
\begin{proof}
  This is the declared model definition of frequency In Hertz.
\end{proof}
\begin{definition}[recorded Answer Choice]
  \label{def:physics:phyx-mini-0295:phyxminiproblems-problemphyxmini0295-recordedanswerchoice}
  \lean{PhyXMiniProblems.ProblemPhyXMini0295.recordedAnswerChoice}
  \uses{def:physics:phyx-mini-0295:phyxminiproblems-problemphyxmini0295-answerchoice}
  The answer label recorded in the supplied dataset metadata
\end{definition}
\begin{proof}
  This is the declared model definition of recorded Answer Choice.
\end{proof}
\begin{definition}[Matches Answer Choice]
  \label{def:physics:phyx-mini-0295:phyxminiproblems-problemphyxmini0295-matchesanswerchoice}
  \lean{PhyXMiniProblems.ProblemPhyXMini0295.MatchesAnswerChoice}
  \uses{def:physics:phyx-mini-0295:phyxminiproblems-problemphyxmini0295-frequencyquantity, def:physics:phyx-mini-0295:phyxminiproblems-problemphyxmini0295-frequencyinhertz, def:physics:phyx-mini-0295:phyxminiproblems-problemphyxmini0295-answerchoice}
  A displayed whole-hertz answer matches the physical frequency when it lies within half a hertz, corresponding to rounding to the nearest hertz
\end{definition}
\begin{proof}
  This is the declared model definition of Matches Answer Choice.
\end{proof}
% --- Archon declaration topology end ---
% --- Archon physics formalization source end ---
